\cotefile{1}
\scanpage{79}

\reversemarginpar
\margin{\begin{tikzpicture}[x=1ex,y=1ex]
\draw[line width=.65pt] (.42,2.4) .. controls (.32,2.15) and (.33,1.95) .. (.53,1.82) -- (.57,.25) .. controls (.57,.05) and (.66,0) .. (.73,.04);
\draw[line width=.65pt] (.45,1.94) .. controls (.19,2.04) and (.12,1.91) .. (.15,1.67);
\end{tikzpicture}}
\normalmarginpar
\uline{Corollaire 4} Soit $u$ un op. de Fredholm dans l'espace de Hilbert $E$. Alors

\[
\text{(18)}\qquad {\color{red}\boxed{\color{black}|\alpha_n(u)|\leq\alpha_n(\deleted{\unsure{\sqrt{u^*u}}}|u|\deleted{{}_{\unsure{2}}})}}
\]

(\unsure{\deleted{En}|En} d'autres termes) les coefficients du développement de Fredholm $\det(1+zu)$ sont majorés en valeur, terme à terme, par les coefficients du déterminant de Fredholm $\det(1+z\unsure{h|H|\lvert u\rvert})$ relatif à l'hermitien positif $h=\sqrt{u^*u}$. On en conclut en effet

\uline{Corollaire 5} Dans les conditions précédentes, on a

\[
\text{(19)}\qquad |\det(1+zu)|\leq\det(1+|z|h)\quad\deleted{\unsure{(\det\sqrt{u^*u})}}
\]

\unsure{Montrons|Prouvons} (18). On a \inserted{en effet}, $|\alpha_nu|=|\mathrm{Tr}\,\overset{n}{\Lambda}u|\leq\|\overset{n}{\Lambda}u\|_1=\alpha_n(\sqrt{u^*u})$ (formule (16)).

\deleted{\uline{Théorème 2} Soit $u$ une \deleted{application}\inserted{op.} de Fredholm dans l'espace de Hilbert $H$. Alors la suite \inserted{$(\lambda_i)$} des valeurs propres \deleted{\illegible{1}} de $u$ est sommable \deleted{\unsure{et}} \deleted{\unsure{(1)|(19)}} \deleted{$\unsure{\sum_i|\lambda_i|\leq\lVert u\rVert_1}$} \deleted{\unsure{valeurs de $u$}}, on a}

\[
\text{\deleted{(20)}}\qquad {\color{red}\boxed{\color{black}\deleted{\det(1+u)=\prod_i(1+\lambda_i)}}}
\]

\deleted{\inserted{Appliquant (21) à $zu$ au lieu de $u$, on trouve}}

\deleted{\uline{Corollaire 1} Le déterminant de Fredholm de $\det(1+zu)$ de $u$ est une fonction entière de genre $0$\,;}

\[
\text{\deleted{(22)}}\qquad \deleted{\det(1+zu)=\prod_i(1+z\lambda_i)}
\]

\reversemarginpar
\margin{Le corollaire 1 ci-\unsure{$=$} équivaut \deleted{\unsure{aux}}\inserted{à la} formule,\par
$\text{(22)}\quad\alpha_n(u)=\sum_{i_1<\mathord{-}<i_n}\lambda_{i_1}\deleted{\unsure{(u)}}\mathord{-}\lambda_{i_n}\deleted{\unsure{(u)}}$\par
dont le plus intéressant est
\begin{tikzpicture}[x=1ex,y=1ex,baseline=0ex]
\draw[line width=.55pt] (.05,1.55) .. controls (.45,1.65) and (.66,1.83) .. (.91,1.99);
\draw[line width=.55pt] (1.45,1.39) -- (.58,.65) -- (.77,.52);
\end{tikzpicture}\par
$\text{(23)}\quad{\color{red}\boxed{\color{black}\mathrm{Tr}\,u=\sum\lambda_i(u)}}$}
\normalmarginpar

\deleted{\uline{Démonstration} du th. 2. \deleted{\unsure{Nous prouvons}} \unsure{On prouve d'abord le lemme} (20) \unsure{sur les valeurs}, \unsure{nous bornant à} la formule (20) : \unsure{soit un} \illegible{4}, \illegible{3} \unsure{propre}, \illegible{5}, \unsure{alors tout revient qu'il en est ainsi}, \unsure{et que} \illegible{1} \unsure{est identique dès que} \deleted{\unsure{tout}} \unsure{la} \illegible{1} \unsure{est nulle}.}

\margin{80}
